% GENERATED, edit the French unit or the translation table
% GENERATED from content/data/enrichment.json by mxbuild.enrichment; do not edit.
\mxunit{enrichment-lattice-paths-notes}{Paths, combinations and Pascal}
\mxsubtitle{High school · counting \textperiodcentered{} Lesson and method \textperiodcentered{} 05}

\begin{prerequisites}
Products, factorials and unordered choices.
\end{prerequisites}

\begin{keypoint}[Ideas and method]
A path from $\pt{0;0}$ to $\pt{a;b}$, using only rightward or upward steps, has $a + b$ steps. Choosing the positions of the a horizontal steps gives $\binom{a+b}{a}$ paths.

For $0 \leq  k \leq  n$, $\binom{n}{k} = \frac{n!}{k!(n-k)!}$. The order of the chosen objects does not matter. Symmetry $\binom{n}{k} = \binom{n}{n-k}$ corresponds to choosing included or excluded objects.

Pascal’s identity $\binom{n}{k} = \binom{n-1}{k-1} + \binom{n-1}{k}$ splits choices by whether one fixed object is included. Paths can also be classified by their last step.
\end{keypoint}

\begin{example}
From $\pt{0;0}$ to $\pt{3;2}$, there are $\binom{5}{3} = 10$ paths. To pass through $\pt{1;1}$, choose one of two paths to that point and one of three paths onward: $2 \times  3 = 6$ paths.
\end{example}

\begin{figurebox}
\begin{tikzpicture}[scale=.85]\draw[mxline](0,0)grid(3,2);\draw[mxmint,very thick,->](0,0)--(1,0)--(1,1)--(2,1)--(3,1)--(3,2);\fill[mxink](1,1)circle(3pt)node[above left]{$\pt{1;1}$};\node[below]at(0,0){$\pt{0;0}$};\node[above]at(3,2){$\pt{3;2}$};\end{tikzpicture}
\end{figurebox}

\begin{activity}[Explore]
Draw a small grid and write the number of arriving paths at each vertex. Add the values of its left and lower neighbours.
\end{activity}

